% Former appendices C (what the corpus shows can be done) and D (mutation detail), removed from the thesis on 5 Nov 2026 at the author's request; kept here as a record. Not built.
\chapter{What the corpus shows can be done}
\label{app:world}

Table~\ref{tab:world} puts the comparison of Section~\ref{sec:res-earlier} in terms of what can be done without this thesis and with it on this corpus; the last row sets it beside Rima et al.'s tool.

\begin{table}[t]\footnotesize
\caption{What the corpus shows can be done without this thesis and with it.}
\label{tab:world}
\begin{tabularx}{\textwidth}{@{}YYY@{}}
\toprule
 & Without this thesis & With this thesis \\
\midrule
Static extraction of the corpus applications & Oshima and Kishi's restrictions admit \EarlierWorkAdmittedApps{}/\EarlierWorkApps{} & extracts all \EarlierWorkApps{} \\
Screens checked without signing in & a browser run renders \ScopeJudged{}/\ScopeScreens{}; the rest redirect to sign-in & every screen the static readers find declared has a node; nothing is run \\
Dead links found from the source & -- & \RoundTwoDeadLinkCandidates{} candidate sites in \RoundTwoDeadLinkWhere{}, checked against the runtime route oracle (Section~\ref{sec:deadlinks}) \\
Transitions on \GTApp{} & its own suite takes \EtoERealworldPairs{} pairs & the graph holds \EtoERealworldGraphPairs{}, an over-approximation \\
Response to a source change & -- & \MutDetected{}/\MutScored{} mutants, at constructs the graph represents \\
Transitions, five corpus applications other than mifos (grouped after reading; relative recall) & Rima et al.'s tool finds \HHCorpusNoMifosTheirsRecall{} of the real pairs either tool finds, at \HHCorpusNoMifosTheirsPrecision{} precision & this extractor finds \HHCorpusNoMifosOursRecall{}, at \HHCorpusNoMifosOursPrecision{}; close on their subjects, unresolved on mifos (Section~\ref{sec:res-h2h}) \\
\bottomrule
\end{tabularx}
\par\smallskip{\footnotesize\noindent\emph{Note.} The dynamic row is this thesis's own rendered-interface check run without signing in, not a crawler with credentials; Oshima and Kishi's tool was not run, Rima et al.'s was (last row, Section~\ref{sec:res-h2h}); every figure is in-sample. All cells are generated figures.}
\end{table}

\chapter{Mutation detail}
\label{app:mutation}

\begin{table}[t]\small
\caption{Detection rate per mutation operator on the attributed graph.}
\label{tab:mutation}
\TabMutation
\par\smallskip{\footnotesize\noindent\emph{Note.} Detected over scored mutants. A dash means the operator found no candidate in that application.}
\end{table}
\FigMutation
